\section*{Chapitre \ref{ch:m2:the-treasury-market} --- Le marché du Trésor}

\begin{solution}{exo:m2:the-treasury-market:1}
27~milliards restent pour les soumissions concurrentielles. 10 + 8 = 18~milliards sont servis à un rendement inférieur
à 4.110~\%~; les 10~milliards à 4.110~\% suffisent pour 9~: taux de coupure 4.110~\%,
attribution au prix de coupure 90.00~\%. \omterm{def:m2:the-treasury-market:tail}{Taux de couverture} (bid-to-cover) $(34 + 3)/30 = 1.23$. Aucun soumissionnaire
n'atteint le cap de 35~\% (10.5~milliards).
\end{solution}

\begin{solution}{exo:m2:the-treasury-market:2}
$4.292 - 4.285 = +0.7$ points de base~: un \omterm{def:m2:the-treasury-market:tail}{tail d'adjudication}. $4.281 - 4.285 = -0.4$~: un stop-through,
le bon résultat, puisque le Trésor a vendu à un rendement plus bas (prix plus élevé)
que celui que payait le marché une minute plus tôt.
\end{solution}

\begin{solution}{exo:m2:the-treasury-market:3}
4.125~\% et 3.875~\%~: les coupons en huitièmes juste en dessous des taux de coupure. Cotée à un
rendement supérieur à son coupon, la note est émise légèrement sous le pair (99.409 pour la
première, une note à dix ans).
\end{solution}

\begin{solution}{exo:m2:the-treasury-market:4}
L'indice de référence est l'émission la plus liquide (fourchette plus serrée, carnets plus profonds)~; elle est
celle avec laquelle tout le monde se couvre et vend à découvert~; elle est donc souvent special en repo et
ses détenteurs la financent à bon compte~; et la demande pour l'émission la plus récente vient des investisseurs
indiciels et de ceux qui recherchent la liquidité. L'opération~: acheter l'\omterm{def:m2:the-treasury-market:otr}{off-the-run},
vendre à découvert l'\omterm{def:m2:the-treasury-market:otr}{on-the-run} pour un même DV01, financer la position longue en repo et emprunter la
position courte par un reverse repo, et attendre que la nouvelle émission devienne ancienne lors de la
prochaine adjudication. Ce qui tourne mal~: le spread s'élargit encore lors d'une fuite vers la
liquidité (mars 2020 l'a fait exactement)~; la position courte devient plus special et le
financement coûte plus cher que ce que rapporte le spread~; l'effet de levier et les appels de marge imposent une sortie
au pire moment.
\end{solution}

\begin{solution}{exo:m2:the-treasury-market:5}
$42\,000/26 = \text{USD}~1\,615$ millions par dealer. La part des dealers, 38.3~\%, correspond
à ce que les investisseurs n'ont pas pris~: une part élevée des dealers signifie que les acheteurs finaux
du marché étaient peu nombreux, et les dealers doivent désormais placer les notes, en général avec une
décote. Le \omterm{def:m2:the-treasury-market:tail}{taux de couverture} (bid-to-cover) compte chaque soumission, y compris celles des dealers
placées très au-dessus du taux de coupure et qui n'avaient aucune chance d'être servies.
\end{solution}

\begin{solution}{exo:m2:the-treasury-market:6}
Le cap représente 35~\% de l'émission~: 14.7~milliards. Il peut acheter les 5.3
milliards restants sur le marché secondaire, via le \omterm{def:m2:the-treasury-market:wi}{marché gris} (when-issued) avant l'adjudication ou après
l'émission, ou auprès des dealers qui se sont vu attribuer plus qu'ils ne voulaient.
\end{solution}

\begin{solution}{exo:m2:the-treasury-market:7}
Le taux de coupure passe de 4.198~\% à 4.200~\%, le \omterm{def:m2:the-treasury-market:tail}{tail d'adjudication} de 1.8 à 2.0 points
de base~: deux dixièmes, et non un point de base entier. Les soumissions indirectes ne représentent
seulement une partie du carnet~; quand elles remontent, les soumissions des dealers et les soumissions directes placées entre
4.198 et 4.208 sont atteintes en premier et absorbent l'essentiel de l'écart. Le taux de coupure
ne suit tout le décalage que si l'ensemble du carnet se déplace.
\end{solution}

\begin{solution}{exo:m2:the-treasury-market:8}
Dans une \omterm{def:m2:the-treasury-market:auction}{adjudication à prix uniforme}, un soumissionnaire retenu paie le taux de coupure, et non sa soumission. Un
petit soumissionnaire dont la soumission ne peut pas déplacer le taux de coupure n'a aucune raison de minorer sa soumission~: il
devrait soumettre le rendement auquel il est indifférent, puisque soumettre un rendement plus élevé
ne fait qu'accroître le risque de ne pas être retenu. Un grand soumissionnaire a un intérêt à agir autrement~: une soumission qui fixe le
taux de coupure modifie le prix payé sur l'ensemble de son attribution, si bien qu'il minore la part de
sa demande qui pourrait être marginale. Dans une adjudication à prix discriminatoire, tout le monde minore,
car chaque soumission est son propre prix.
\end{solution}

\begin{solution}{pb:m2:the-treasury-market:1}
\textbf{1.} $42\,000 - 300 = \text{USD}~41\,700$ millions.
\textbf{2.} Taux de coupure 4.198~\%~; \omterm{def:m2:the-treasury-market:tail}{tail d'adjudication} $+1.8$ points de base.
\textbf{3.} \omterm{def:m2:the-treasury-market:tail}{Taux de couverture} (bid-to-cover) 1.82~; attribution au prix de coupure 95.02~\%.
\textbf{4.} Soumissions indirectes 51.1~\%, directes 10.6~\%, dealers 38.3~\%.
\textbf{5.} Coupon 4.125~\%~; prix 99.409 pour 100.
\textbf{6.} $42\,000/26 = 1\,615.4$ millions~; il soumet 1\,615 millions à trois
rendements, ce qui respecte l'obligation si les rendements sont raisonnablement compétitifs.
\textbf{7.} Les soumissions à 4.188~\% et 4.195~\% sont sous le taux de coupure et servies
en totalité~; la soumission à 4.205~\% est au-dessus et n'obtient rien~: attribution de 1\,000
millions.
\textbf{8.} $807 \times 1\,000 = \text{USD}~806\,797$ par point de base.
\textbf{9.} Elle a vendu à 4.180~\% et acheté à 4.198~\%~: 1.8 points de base en sa
faveur, moins 0.25 de coûts, sur 806\,797 par point de base~: $+\text{USD}~1\,250\,536$.
\textbf{10.} Un stop-through de 0.2 points de base~: $808 \times 1\,000 \times (-0.2
- 0.25) = -\text{USD}~363\,758$.
\textbf{11.} Avec un \omterm{def:m2:the-treasury-market:tail}{tail d'adjudication} de $+0.25$ points de base, le coût.
\textbf{12.} Les acheteurs finaux ont soumis moins que ce qui était offert au prix du marché,
si bien que le Trésor a dû vendre plus loin dans les soumissions de repli des dealers~;
la note du \omterm{def:m2:the-treasury-market:wi}{marché gris}, devenue la note émise, se déprécie vers le taux de coupure à mesure que le
marché absorbe les notes restées chez les dealers.
\textbf{13.} Pour respecter l'obligation au prorata, et parce qu'une soumission située à quelques dixièmes de
point de base au-dessus du marché est le prix de la prise ferme~: elle est retenue si
l'adjudication est faible, et elle est alors rémunérée pour le risque.
\textbf{14.} À 4.200~\%, un \omterm{def:m2:the-treasury-market:tail}{tail d'adjudication} de 2.0 points de base (l'\cref{exo:m2:the-treasury-market:7}).
\textbf{15.} Il aurait soumis moins et vendu moins sur le \omterm{def:m2:the-treasury-market:wi}{marché gris} (when-issued), ou soumis un rendement plus bas
pour n'être servi que dans une adjudication faible. Il ne peut pas connaître le \omterm{def:m2:the-treasury-market:tail}{tail d'adjudication}~: le
niveau du \omterm{def:m2:the-treasury-market:wi}{marché gris} (when-issued) est la meilleure estimation du taux de coupure par le marché, et le P\&L du dealer
est un pari sur l'écart.
\textbf{16.} Un prix uniforme supprime pour la plupart des soumissionnaires la raison de minorer, attire davantage
de soumissionnaires (les petits n'ont pas à deviner les soumissions des autres), et fait que chaque
soumissionnaire retenu paie le même prix, celui auquel le marché s'est équilibré. Le Trésor l'a utilisé
pour toutes ses adjudications depuis novembre 1998.
\textbf{17.} Certains ne peuvent pas soumissionner directement ou ne veulent pas révéler leur demande~;
acheter sur le \omterm{def:m2:the-treasury-market:wi}{marché gris} (when-issued) auprès d'un dealer fixe leur prix à l'avance, et le dealer
prend le risque de l'adjudication à leur place.
\textbf{18.} Dans une adjudication faible, il est servi sur des soumissions qu'il aurait préféré perdre, et
il porte les notes à perte jusqu'à leur placement. En contrepartie~: le statut de
\omterm{def:m2:the-treasury-market:pd}{spécialiste en valeurs du Trésor} (primary dealer), une relation de contrepartie avec la banque centrale, et les
flux clients.
\textbf{19.} Résultat à retenir~: \emph{le P\&L de tail du dealer}~: USD~806\,797 par
point de base de tail sur son attribution de USD~1 milliard~; USD~1\,250\,536 dans cette adjudication
(tail de 1.8 points de base)~; seuil de rentabilité à un tail de 0.25 point de base.
\textbf{20.} L'écart entre ce que l'adjudication a payé et ce que le marché
payait à la clôture des soumissions.
\end{solution}

\begin{solution}{iq:m2:the-treasury-market:1}
Le rendement de coupure moins le rendement du \omterm{def:m2:the-treasury-market:wi}{marché gris} à la date limite de soumission. Un \omterm{def:m2:the-treasury-market:tail}{tail d'adjudication} positif et important
signifie que l'adjudication a dû monter vers des rendements plus élevés que ceux auxquels le marché échangeait, pour être servie~:
la demande finale était faible, les dealers ont pris davantage que d'habitude,
et le marché se déprécie généralement après le résultat.
\iqlookfor{la définition avec son signe, et le lien avec la prise ferme par le dealer.}
\end{solution}

\begin{solution}{iq:m2:the-treasury-market:2}
La liquidité (le titre de référence est celui où se concentrent les négociations, la couverture et la livraison des contrats à terme),
la specialness en repo (c'est l'émission la plus vendue à découvert, donc la détenir permet d'emprunter du cash à bon marché),
et la demande des investisseurs qui ont besoin du titre de référence.
Les détenteurs acceptent de renoncer à un peu de rendement pour ces avantages~; le spread varie avec la prime de liquidité
et s'élargit en période de stress.
\iqlookfor{la liquidité et le financement, et non seulement «~il est plus récent~».}
\end{solution}

\begin{solution}{iq:m2:the-treasury-market:3}
Dans une adjudication à prix multiples, chaque gagnant paie sa propre soumission, si bien que les soumissionnaires
ajustent leurs soumissions vers le rendement de coupure attendu et doivent deviner les soumissions des autres~; les gagnants le regrettent
(malédiction du vainqueur). Dans une \omterm{def:m2:the-treasury-market:auction}{adjudication à prix uniforme}, les gagnants paient le rendement d'équilibre,
si bien que les petits soumissionnaires offrent leur rendement de réserve~; seuls les plus gros, capables de fixer le prix,
ont intérêt à ajuster leurs soumissions.
\iqlookfor{qui paie quoi, et quels soumissionnaires continuent d'ajuster leurs soumissions.}
\end{solution}

\begin{solution}{iq:m2:the-treasury-market:4}
Après une publication économique décevante sur les ventes au détail, la profondeur des carnets d'ordres à dix ans est tombée à une
fraction de sa valeur normale~; entre 9:33 et 9:45, les rendements ont baissé de 16 points de base et
sont revenus sur leurs niveaux sans nouvelle information, alors que les sociétés de trading pour compte propre ont réalisé plus de la moitié
du volume. Dans les données des carnets d'ordres~: la profondeur aux meilleurs niveaux et à quelques ticks du meilleur cours,
au fil du temps, les annulations et les ratios ordres sur exécutions, la part du volume
par type de participant, les auto-transactions, le volume agressif face au volume passif par
participant, et l'avance-retard entre contrats à terme et comptant.
\iqlookfor{la profondeur, et non le volume, comme variable, et une liste concrète de
grandeurs mesurables.}
\end{solution}

\begin{solution}{iq:m2:the-treasury-market:5}
Parce que tout le monde avait besoin de liquidités, et que les \omterm{def:m2:money-markets:bill}{bons du Trésor} sont ce qui peut être vendu~: les investisseurs à effet de levier dénouaient leurs positions,
les détenteurs étrangers levaient des dollars et les fonds confrontés à des rachats les ont vendus,
et les dealers ne pouvaient pas absorber le flux dans leurs bilans.
L'actif sans risque était liquide, donc il a été vendu en premier~; les achats de la banque centrale
plus de mille milliards de dollars en cinq semaines ont rétabli le fonctionnement du marché.
\iqlookfor{la distinction entre sécurité et liquidité, et la capacité
d'intermédiation des dealers.}
\end{solution}

\begin{solution}{iq:m2:the-treasury-market:6}
Entités~: titre (CUSIP, durée, dates d'émission et d'échéance), annonce
(émission, dates), cours du \omterm{def:m2:the-treasury-market:wi}{marché gris} (horodatés, avec leur source), résultat (rendement
de coupure, attribution au taux le plus élevé, \omterm{def:m2:the-treasury-market:tail}{taux de couverture}, parts des soumissionnaires, coupon, prix),
tail calculé. Chemin en temps réel~: s'abonner à la publication du résultat, analyser, chercher
le cours du \omterm{def:m2:the-treasury-market:wi}{marché gris} à la date limite dans un historique indexé par le temps, calculer,
publier. Tests~: analyser les résultats publiés d'adjudications passées et les comparer à leurs statistiques publiées~;
un cours du \omterm{def:m2:the-treasury-market:wi}{marché gris} exactement à la date limite, et les cours manquants~;
les émissions réouvertes (coupon déjà connu, cotées en prix)~;
les bons (taux d'escompte, et non rendements)~; la gestion de l'horloge et du fuseau horaire
de la date limite.
\iqlookfor{l'horodatage de la date limite comme point central, et les tests
sur l'historique publié.}
\end{solution}
